\documentclass[11pt]{article}
\usepackage[letterpaper,margin=0.85in]{geometry}
\usepackage[T1]{fontenc}
\usepackage{lmodern,amsmath,booktabs,array,xcolor,fancyhdr}
\usepackage[colorlinks=true,linkcolor=blue,urlcolor=blue]{hyperref}
\definecolor{ink}{HTML}{17344A}
\pagestyle{fancy}\fancyhf{}
\fancyhead[L]{\small\textcolor{ink}{UpstreamDrift / Necromatcher}}
\fancyhead[R]{\small Replay Impact Application Methods}
\fancyfoot[L]{\small Authored Research; Historical Qualification Pending}
\fancyfoot[R]{\thepage}
\setlength{\headheight}{14pt}
\setlength{\parindent}{0pt}\setlength{\parskip}{7pt}
\newcommand{\code}[1]{\nolinkurl{#1}}
\title{\textcolor{ink}{Necromatcher Replay Impact Application\\Execution, Recall and Repeatability}}
\author{UpstreamDrift Technical Reference\\Epic \#11232; Integration Issue \#11464}
\date{4 October 2026 / Application Supplement 1}
\begin{document}
\maketitle
\begin{center}\rule{0.9\linewidth}{0.6pt}\end{center}
\section*{Purpose and Scientific Boundary}
This supplement documents the application journey from a registered authored
replay to an owned native impact job, retained trajectory and portable state
receipt. It complements the five-page extraction methods reference and the
camera/source-scope report; their earlier evidence and producer identities remain
unchanged. Exact software, test, toolchain and artifact hashes belong to the
accompanying application review receipt. A source fingerprint includes working
implementation bytes; a recorded Git commit alone does not identify dirty code.

Recorded rates refer to authored simulation seconds. Geometry, effective mass
and inertia, the selected sample and the flight transform are operator
assumptions. Both scientific and physical-source-time qualification remain
false. A successful job is a completed software calculation with research
acceptance \emph{rejected}; it is not a qualified historical impact prediction.

Tiger retains original frames 0--190, the reviewed half-open interval $[0,191)$.
Frames 191 onward, including released-hand follow-through, are excluded. The
cutoff is conservative; exact physical release and contact remain unmeasured.
Hogan retains $[0,750)$. No application or impact success changes these domains.

\section*{Ownership Chain}
\begin{center}\small
\begin{tabular}{>{\raggedright\arraybackslash}p{0.25\linewidth}>{\raggedright\arraybackslash}p{0.66\linewidth}}\toprule
Boundary & Responsibility\\\midrule
React Controls & Verified replay recall, explicit declarations, budget, job recall.\\
Local API & Typed admission, replay ownership, bounded job dispatch and download.\\
Matching Job Service & Persistent state, cancellation and terminal publication.\\
Native Worker & Recorded-state extraction, existing impact/flight calculation.\\
Public Receipt Owner & Exact trajectory-byte binding and detached state recall.\\
Existing Viewers & Retained sample import without re-simulation.\\\bottomrule
\end{tabular}
\end{center}
No independent forward kinematics, collision solver or trajectory renderer is
introduced by this application boundary.

\newpage
\section*{Explicit Request and Admission}
The operator imports a finite JSON object containing exactly \code{geometry}
and \code{selection}. Geometry declares a native body, body-local head point
in metres, unit perpendicular face normal/up vectors, positive effective mass
in kg and scalar inertia in kg\,m$^2$, and an assumption description. Selection
declares an existing recorded sample, a proper orthonormal world-to-flight
rotation, translation in metres and a selection description. The application
requires an explicit positive execution budget no greater than 600 seconds.
It does not manufacture a contact sample, normalize malformed axes or infer
physical rates from video timestamps.

Server admission delegates to the existing typed declaration owners and
canonical Library authentication. The replay, torque profile, fit, native model
and capture must retain their expected kinds, session and exact hashes. A sample
outside the stored trace, foreign replay identity, malformed declaration or
unauthenticated parent is refused before scheduling native work. The local-client
restriction remains in force. Client requests contain no output directory.

\section*{Recorded Kinematics and Flight Frame}
The public native marker linearizer supplies the point Jacobian and face-triad
derivatives at the recorded coordinates. With recorded generalized rate $v$,
\[
\dot p=J(q)v,\qquad
\omega=\tfrac12\sum_{i=1}^{3}e_i\times\dot e_i.
\]
The extraction methods reference derives this rigid-triad identity and its
orthonormality/tangent checks. For the authored proper rigid map $(R,t)$,
\[
p_f=Rp+t,\quad v_f=R\dot p,\quad \omega_f=R\omega,\quad n_f=Re_1.
\]
Translation affects positions only. The existing consumer preserves the full
velocity, angular velocity and face normal declared in
\code{flight_xfwd_yleft_zup}; it does not replace them with speed and nominal loft.

\section*{Execution Identity}
An impact-specific stamp extends the canonical refit stamp with impact/flight
Python sources and the two canonical trajectory validators. Runtime identity
also retains the installed flight distribution's version and executable-file
hashes, detecting same-version binary replacement. Three stable fields are
rechecked: source commit, source-map digest and runtime digest. Start timestamps
are retained as provenance but are not compared for equality across calls.
Missing flight capability fails through the existing pipeline; production
execution does not substitute a test trajectory or repair the environment.

\newpage
\section*{Bounded Worker and Publication}
The server creates a UUID-named run beneath the Library's \code{impact-runs}
directory. The fixed clean-process transport dispatches the reviewed impact
worker, imports MuJoCo before its execution-specific native APIs and bounds the
owned process by cancellation and the declared wall budget. It accepts no
arbitrary worker module or user-selected publication path.

The worker uses public replay extraction, the existing trajectory handoff
coordinator and actual impact/flight pipeline. Parent and stable execution
identities are rechecked around calculation. Output is staged in an owned
temporary directory and published only as a complete bundle to an absent
destination. Existing or concurrently appearing output is refused. Cancellation,
worker failure, stale parents or changed source/runtime cannot publish a verified
completion record. The matching service retains failed/cancelled state.

The job owner verifies canonical receipt contents, source assumptions and all
artifact hashes before terminal publication. A restarted host may recall
completed artifacts but cannot invent a live cancellation handle for an orphan
pending/running record. Download repeats parent and artifact validation before
and after ZIP creation; incomplete or changed artifacts are refused.

\section*{Portable Bundle and Effective Settings}
\begin{center}\small
\begin{tabular}{>{\raggedright\arraybackslash}p{0.29\linewidth}>{\raggedright\arraybackslash}p{0.62\linewidth}}\toprule
Artifact & Retained Content\\\midrule
\code{trajectory.json} & Canonical six-field trajectory wire and retained samples.\\
\code{impact-receipt.json} & Versioned envelope: state, exact trajectory SHA-256 and size.\\
\code{result.json} & Parent chain, declarations, execution stamp, calculation settings/results.\\
\code{request.json} & Server-owned run identity, declared budget and admitted assumptions.\\\bottomrule
\end{tabular}
\end{center}
The result retains effective environment and impact parameters, ball assumptions,
launch conditions, post-impact state, trajectory metrics and full extraction
metadata. Default method/flight settings are explicit assumptions of the pinned
pipeline implementation; they are not historical calibrations. This action does
not expose arbitrary environmental or impact-model configuration.

The public receipt loader binds exact trajectory bytes, validates the canonical
viewer wire and returns detached state vectors without SDK execution, impact
calculation or resampling. Portable parent hashes remain provenance claims;
receipt-only recall is not fresh Library authentication. Job recall adds that
Library authentication separately. The trajectory alone carries no complete
historical receipt, so retain and distribute the bundle together.

\newpage
\section*{Repeatable Operator Procedure}
\begin{enumerate}
\item Open Necromatcher, select the player/swing and recall a registered authored replay.
Inspect its parent hashes, authored clock and unqualified status.
\item Import exact geometry/selection declarations and enter a bounded wall budget.
Review the native body, local point, face axes, effective inertia and sample.
\item Choose \emph{Preview Research Impact}. Observe status, rejection blockers and
the run ID. Cancel through the owned job when required.
\item After verified completion, inspect the saved sample, declarations and metrics.
They belong to the result, even if a different declaration is now loaded in the form.
\item Download the verified four-file bundle. Retain its run ID. Recall that ID
through \emph{Recall Research Impact Run}; reload does not require re-running.
\item Extract \code{trajectory.json} and use the existing Ball Flight viewer's
trajectory import. Retained samples are imported; do not treat a new simulation
as recall. Preserve the companion receipt when interpreting or sharing it.
\end{enumerate}

\section*{Verification and Repeatability Evidence}
Meaningful failing tests cover absent integration, missing stored-run recall,
false display of unverified artifacts and timestamp comparison that incorrectly
rejected otherwise identical execution stamps. Unit tests cover strict geometry,
sample bounds, parent/source/runtime changes, output ownership, rollback,
cancellation, foreign identity, incomplete publication and byte tampering.

Native acceptance uses a temporary canonical Library and synthetic model/profile.
It executes independent MuJoCo replay, registered HDF5 recall, the clean impact
worker and canonical Rust flight. Native extraction, solvers, stamps and transport
are not substituted. Checks cover portable state, viewer imports, restarted recall
and trajectory-byte tampering. This is software evidence, not a historical capture.
The review receipt pins closed test logs, sources and the local flight-wheel build.

\section*{Remaining Historical Acceptance}
Both historical fits still require improved reprojection and justified camera,
anthropometry, range-of-motion, contact and physical-time evidence. Rolling
shutter remains a conditional sensitivity model; shaft flex, blur, interlacing,
film scanning and broadcast processing remain alternative explanations.
The full goal also requires historical authored dynamics and integration into
simulation golf. A successful synthetic application handoff does not complete
those requirements or establish engineering-manual release approval.

\section*{References and Public Owners}
\href{https://github.com/D-sorganization/UpstreamDrift/issues/11464}{Issue \#11464}
and \href{https://github.com/D-sorganization/UpstreamDrift/pull/11359}{PR \#11359}.
Public owners are exported by the workspace facade. The repeatable Markdown
procedure lists their interfaces. Earlier extraction and camera/source-scope
reports retain their original evidence boundaries.
\end{document}
